\section*{Chapter 1: Introduction}
\setcounter{section}{1} \setcounter{subsection}{0}
\phantomsection \addcontentsline{toc}{section}{\textit{Chapter 1: Introduction}}

In recent years, LLMs have become the cornerstone of conversational AI, enabling a wide range of AI agents \cite{park2023generative, qin2024toolllm}. 
Much like the steam engine in the Second Industrial Revolution, LLMs serve as a general-purpose driving force, powering the emergence of diverse agentic applications.
Recent advances in LLM capabilities have enabled instruction following, reasoning, and interaction with external environments, forming the foundation for agentic systems \cite{yao2023react, wang2024agentbench, schick2023toolformer}.
As a result, emergent AI agents can present strong task fulfillment abilities in diverse environments, ranging from direct acting assistants to more sophisticated reasoning and problem-solving systems \cite{mei2024aios,liu2026agentos}.

\subsection{Problem Statement}

Tool-using, stateful AI agents expose a fundamental mismatch between modern agent frameworks and the abstractions provided by conventional operating systems \cite{zhang2024llmos, wu2024agentsystems}. 
While the existing systems effectively manage processes, files, and network resources, MCP-style workflows are organized around higher-level concepts such as semantic tool descriptions, session context, approval boundaries, and cross-step state.
This mismatch leads to fragile control planes when enforcement remains entirely in user space. 
Frameworks such as OpenClaw\cite{openclaw2026,openclaw_security1} demonstrate that granting agents execution capability and persistent context can lead to security risks including prompt injection, privilege escalation, and cross-step exploitation. 

A single gateway process is typically responsible for interpreting tool manifests, binding requests to identities, and authorizing execution, effectively placing the control plane in user space. 
Such a design introduces a single point of failure and expands the trusted computing base, making it vulnerable to crashes, logic bugs, or compromised control paths~\cite{baumann2014arrakis}.
As a result, critical safety properties become difficult to guarantee, including ensuring that executed tools match approved semantics, that requests remain bound to a trusted identity, and that approval state persists across failures \cite{greshake2023prompt, qin2024toolllm}.

These observations motivate a key question: which control plane invariants in tool invocation fundamentally benefit from kernel-level enforcement, and can such enforcement improve security and observability without sacrificing performance?

\subsection{Our Solution}

To mitigate the above-mentioned problems, we introduce \texttt{linux-mcp}, which isolates a smaller kernel responsibility: enforcing control-plane invariants for tool invocation. 
Rather than proposing a full AI-native operating system, \texttt{linux-mcp} keeps semantic parsing and tool execution in user space while moving admission control, approval state, binding checks, and durable observability into a Linux kernel module.

The design follows a simple placement principle: the kernel holds the control plane, while user space retains the execution plane. 
This avoids pushing complex semantics or application logic into the kernel, while still allowing critical control decisions to be enforced outside the gateway process.

The architecture has four layers: an LLM client (\texttt{llm-app}), a manifest-aware gateway (\texttt{mcpd}), a Generic-Netlink-based kernel control plane (\texttt{kernel\_mcp}), and resident tool services over Unix domain sockets. 
The \texttt{mcpd} gateway loads manifests, computes semantic identifiers, validates requests against schemas, binds sessions to Unix-domain peer credentials, and forwards tool invocations through the kernel mediation path. 
The kernel module maintains compact tool and agent registries, performs arbitration based on binding and semantic metadata, and manages approval tickets for high-risk operations. 
Tool services execute concrete operations behind Unix domain sockets and are not directly accessible to the client.

Requests follow a fixed path: \texttt{llm-app} issues a tool invocation to \texttt{mcpd}, which validates and forwards it to the kernel; the kernel returns an allow, deny, or defer decision; allowed requests are executed by tool services, and completion is reported back through the same path. 
This keeps all control-plane transitions explicit and observable.

By making these invariants kernel-visible, the system can enforce them with higher integrity and improved observability.
Semantic consistency is enforced through manifest-derived identifiers; identity binding is anchored in OS-level peer credentials; approval state is maintained in the kernel across daemon failures.
All tool invocations pass through a single mediation path, reducing reliance on transient in-memory gateway state.

Critically, this design does not reduce the trusted footprint of \texttt{mcpd}---the gateway remains fully trusted and populates the kernel's registries.
What it achieves is a \emph{partitioning} of control-plane state: registry entries, approval tickets, and binding checks are held in the kernel under \texttt{CAP\_NET\_ADMIN}-protected Netlink access control, making them durable across daemon failure and inaccessible to unprivileged processes, while semantic parsing and tool execution remain in user space where flexibility matters most.

This design does not move tool execution into the kernel or attempt to replace existing Linux security mechanisms. 
Instead, it introduces a kernel-assisted control plane that complements user-space flexibility while strengthening enforcement for MCP-style tool invocation.

\paragraph{Comparison with user-space systems.}
To highlight the key design differences, we compare \texttt{linux-mcp} with a representative user-space system, OpenClaw.

\begin{table}[t]
\centering
\small
\begin{tabular}{lcc}
\toprule
\textbf{Feature} & \textbf{linux-mcp} & \textbf{OpenClaw} \\
\midrule

Control-plane placement 
& Kernel-assisted 
& User-space gateway \\

Trust boundary 
& Split (user + kernel) 
& Monolithic (user-space only) \\

Approval enforcement 
& Kernel-enforced (ticket-based) 
& User-space logic \\

Binding correctness 
& Strong (kernel-verified) 
& Best-effort (gateway-managed) \\

Failure resilience 
& State persists in kernel 
& Lost on gateway failure \\

Security model 
& Minimal TCB (kernel invariants) 
& Large TCB (gateway-dependent) \\

\bottomrule
\end{tabular}
\caption{Key design differences between \texttt{linux-mcp} and OpenClaw.}
\label{tab:comparison}
\end{table}

\noindent
\textbf{Summary.}
\texttt{linux-mcp} moves critical control-plane invariants into the kernel, reducing reliance on user-space correctness while improving binding safety and failure resilience.




\subsection{Main Contributions}

We build a runnable end-to-end prototype, \texttt{linux-mcp}, that realizes a kernel-assisted control plane for tool invocation on Linux. 
The system combines a Linux kernel module for maintaining tool and agent registries, enforcing mediation decisions, and exposing state via \texttt{sysfs}, with a user-space gateway (\texttt{mcpd}) that handles manifest loading, session binding, and request forwarding, and a set of tool services accessed through a unified mediation path.

Our design deliberately avoids moving tool execution or semantic parsing into the kernel. 
Instead, \texttt{linux-mcp} focuses on kernel-assisted control-plane enforcement, preserving user-space flexibility while strengthening security guarantees.

This paper makes the following contributions:
\begin{itemize}

    \item We design and implement \texttt{linux-mcp}, a kernel-assisted architecture that partitions control-plane state between kernel and user space: the kernel holds registry entries, approval tickets, and binding checks under \texttt{CAP\_NET\_ADMIN}-enforced Netlink access control, while semantic parsing and tool execution remain in user space. This partitioning makes critical state durable across daemon failure and inaccessible to unprivileged processes, without reducing the trusted footprint of \texttt{mcpd}.

    \item We formulate kernel-assisted control of tool invocation as a systems problem for AI agent workloads on Linux, identifying a mismatch between existing operating system abstractions and agent-centric execution models, and precisely characterising which control-plane invariants benefit from kernel placement.

    \item We conduct a comprehensive evaluation against pure and hardened user-space baselines across latency, throughput, attack resistance, and failure scenarios, showing that \texttt{linux-mcp} eliminates all tested control-plane attack classes with negligible overhead, and that kernel-held approval state survives daemon failure in a way neither baseline can match.
\end{itemize}
